\documentclass[11pt]{article}

\usepackage[margin=1in]{geometry}
\usepackage{amsmath,amssymb,amsthm}
\usepackage[hidelinks]{hyperref}
\urlstyle{same}
\emergencystretch=2em

\newtheorem{theorem}{Theorem}[section]
\newtheorem{proposition}[theorem]{Proposition}
\newtheorem{lemma}[theorem]{Lemma}
\newtheorem{corollary}[theorem]{Corollary}
\newtheorem*{thmA}{Theorem A}
\theoremstyle{definition}
\newtheorem{definition}[theorem]{Definition}
\theoremstyle{remark}
\newtheorem{remark}[theorem]{Remark}

\newcommand{\On}{\mathbf{On}}
\newcommand{\Nat}{\mathbb N}
\newcommand{\Integers}{\mathbb Z}
\newcommand{\Rationals}{\mathbb Q}
\newcommand{\Hess}{\mathcal H}
\newcommand{\Pair}{\mathcal P}
\newcommand{\Connected}{\mathcal C}
\newcommand{\gr}{\operatorname{gr}}
\newcommand{\rank}{\operatorname{rank}}
\newcommand{\Span}{\operatorname{span}}

\title{Stable Quadratic Pairs over Unique-Base Semirings:\\
Hessenberg Collapse, Supertropical Blocks, and a Wall Obstruction}
\author{GPT 5.6 Sol \and a9lim}
\date{}

\begin{document}
\maketitle

\begin{abstract}
We construct a stable theory of finite free quadratic pairs over semirings
whose free modules have unique bases.  Such a pair is a finite weighted graph:
vertices carry the quadratic and companion-diagonal coefficients, edges carry
the off-diagonal companion coefficients, monomial basis changes act by a unit
gauge, and orthogonal sum is disjoint union.  The isometry monoid is therefore
free commutative on connected gauge-orbit types.  Regular pairs split into
rank-one loops and rank-two edges; basic metabolic stabilization deletes
exactly the metabolic edge generators.

For the closed Hessenberg fragment
\(\Hess_d=\{\alpha<\omega^{\omega^d}\}\), natural sum and product identify
\(\Hess_d\) with \(\Nat[x_0,\ldots,x_{d-1}]\).  Its only unit is one and
additive cancellation makes the companion unique and balanced.  Every regular
pair is consequently a sum of copies of one hyperbolic plane, so the regular
Witt group is zero.  The all-pairs stable monoid remains nontrivial: it is
free on all
connected weighted graphs except the hyperbolic edge, and has an exact finite
decision procedure.

For the standard rational tangible supertropical cover we give the analogous
gauge-orbit presentation, the complete rank-two companion ambiguity, and the
regular metabolic generators.  Coefficientwise supertropicalization is
functorial only for based forms.  We compute the universal further collapse:
for every field of characteristic different from two and every rational
tangible supervaluation, the basis-independent ring-image quotient is exactly
\(\mathbb Z/2\), detected by rank parity, for both the canonical balanced and
polar-companion lifts.  No other classical Witt datum survives.

Finally, completing thermograph walls by one global tropical zero produces
honest max-plus and min-plus coefficient semifields.  Each wall pair defines a
regular rank-one quadratic pair classified exactly by its wall gap.  The hot
switch \(\{1\mid-1\}\) is self-negative but has nonzero gap, proving that the
resulting stable wall class cannot be additive under disjunctive sum.
\end{abstract}

\section{Introduction}

Over a ring, subtraction hides three different operations in the usual Witt
construction: recovery of the polar form from the quadratic form, orthogonal
cancellation, and passage from a monoid to a group.  None is automatic over a
semiring.  The correct first object is therefore a quadratic pair \((q,b)\),
where
\[
 q(ax)=a^2q(x),\qquad
 q(x+y)=q(x)+q(y)+b(x,y),
\]
and \(b\) is a specified symmetric bilinear companion
\cite{ikr2016qf1}.  Isometries in this paper preserve both entries of the
pair.

The unique-base theorem of Izhakian, Knebusch, and Rowen is the organizing
input: if the nonzero elements of a semiring are closed under addition and
multiplication, every basis of a free module differs from every other by a
permutation and multiplication by scalar units
\cite[Theorem~1.3]{ikr2016uniquebase}.  This replaces matrix reduction by
finite graph theory.  It also makes clear why the regular theory and the
singular envelope behave differently.

The main conclusions are as follows.

\begin{thmA}
Let stable equivalence kill finite free regular pairs carrying a basic
Lagrangian, as in Definition~\ref{def:stable}.
\begin{enumerate}
\item For every \(d\ge0\), all regular quadratic pairs over the Hessenberg
fragment \(\Hess_d\) are hyperbolic.  Its regular Witt group is zero.  The
stable monoid of all pairs is the free commutative monoid on connected
weighted-graph types other than the hyperbolic plane.
\item Over the standard rational tangible supertropical cover, the all-pairs
stable group is free abelian on connected coefficient-graph gauge orbits
other than the regular metabolic edge orbits.  Regular pairs have only loop
and edge blocks, and rank-two companion ambiguity is explicit.
\item Canonical coefficientwise supertropicalization neither descends from
based forms to unbased isometry classes nor preserves regularity under an
arbitrary classical change of basis.  For a source field of characteristic
different from two, imposing all classical basis relations makes its
universal nonsingular ring-image quotient exactly
\(\mathbb Z/2\), detected by rank parity.  The same is true for the alternative
polar-companion lift.
\item The max/min thermograph walls do define regular rank-one stable classes
over completed piecewise-linear tropical semifields.  Their class is the wall
gap, but it cannot be additive in the Conway game group.
\end{enumerate}
\end{thmA}

The Hessenberg arm is a complete classification and collapse theorem.  The
supertropical block theorem and the universal classical ring-image quotient
are complete for the stated coefficient cover.  The wall theorem is a no-go
theorem for the natural max/min construction, not a claim that every possible
synchronized cross-quotient is impossible.

\section{The unique-base stable monoid}

Let \(R\) be a commutative semiring such that \(R\setminus\{0\}\) is closed
under addition and multiplication.  Let \(V=R^n\) with basis
\(E=(e_1,\ldots,e_n)\), and let \(P=(V,q,b)\) be a quadratic pair.

\begin{definition}[Coefficient graph]\label{def:graph}
The coefficient graph \(G_E(P)\) has vertex set \(\{1,\ldots,n\}\).  Vertex
\(i\) is labelled by
\[
 (a_i,d_i)=(q(e_i),b(e_i,e_i)).
\]
For \(i<j\), there is an edge \(ij\) precisely when
\(c_{ij}=b(e_i,e_j)\ne0\), and that edge is labelled by \(c_{ij}\).
Only labelings that occur in a quadratic pair are called admissible.
\end{definition}

If \(u_i\in R^\times\) and \(\pi\in S_n\), the monomial basis change
\(e_i\mapsto u_i e_{\pi(i)}\) acts by
\[
 a_i\mapsto u_i^2a_{\pi(i)},\qquad
 d_i\mapsto u_i^2d_{\pi(i)},\qquad
 c_{ij}\mapsto u_iu_jc_{\pi(i)\pi(j)}.
\]
We call this the unit gauge action.

\begin{proposition}[Coordinate expansion]\label{prop:expansion}
For \(x=\sum_i x_ie_i\),
\[
 q(x)=\sum_i x_i^2a_i+\sum_{i<j}x_ix_jc_{ij}.
\]
Thus \(G_E(P)\), including the diagonal labels \(d_i\), determines the pair.
\end{proposition}

\begin{proof}
Apply the second quadratic-pair identity repeatedly.  Bilinearity gives
\(b(x_ie_i,x_je_j)=x_ix_jc_{ij}\), and homogeneity gives
\(q(x_ie_i)=x_i^2a_i\).  A bilinear form is determined by its matrix on the
basis, so the additional labels \(d_i\) determine the full companion.
\end{proof}

\begin{theorem}[Graph presentation]\label{thm:graph}
Let \(\Connected_R\) be the set of connected admissible coefficient graphs
modulo the unit gauge action and labelled-graph isomorphism.  The monoid of
isometry classes of finite free quadratic pairs satisfies
\[
 \Pair(R)\cong\Nat^{(\Connected_R)}.
\]
In particular, orthogonal cancellation holds.
\end{theorem}

\begin{proof}
The unique-base theorem makes every module automorphism monomial.  By
Proposition~\ref{prop:expansion}, such a map is an isometry of pairs exactly
when its permutation and units obey the displayed gauge law.  An orthogonal
decomposition along subsets of the basis is equivalent to the absence of
cross edges.  Hence the indecomposable orthogonal summands are exactly the
connected components of the coefficient graph.  Finite labelled graphs have
unique connected-component multisets, which gives the displayed free
commutative monoid and cancellation.  This is the pair-valued specialization
of the component and cancellation theory in
\cite[Theorems~2.6 and~4.9]{ikr2016uniquebase}.
\end{proof}

Write \(b^\flat:V\to V^*=\operatorname{Hom}_R(V,R)\) for
\(b^\flat(x)=b(x,-)\).

\begin{definition}[Regular, split, and metabolic]\label{def:stable}
A pair is \emph{regular} if \(b^\flat\) is an isomorphism.  A submodule is
\emph{basic} if it is spanned by a subset of the unique basis.  A regular pair
is \emph{basic metabolic} if it has a basic submodule \(L\) such that
\(q|_L=0\) and \(L=L^\perp\).  It is \emph{basic split} if it has complementary
basic Lagrangians.

Let \(\mathcal N_R\) be the monoid of basic metabolic pairs.  Define
\[
 A\sim_{\mathrm{st}}B
 \quad\Longleftrightarrow\quad
 \exists N_1,N_2\in\mathcal N_R:\
 A\perp N_1\cong B\perp N_2.
\]
The all-pairs stable group is
\[
 W_{\mathrm{all}}(R)=K_0(\Pair(R))/\langle\mathcal N_R\rangle.
\]
The regular Witt group is defined by the same construction inside the regular
submonoid.
\end{definition}

Allowing different neutral objects on the two sides is essential.  Adding the
same object to both sides gives no new equivalence because Theorem~\ref{thm:graph}
already proves cancellation.

\begin{theorem}[Regular block theorem]\label{thm:blocks}
Every regular quadratic pair over \(R\) is an orthogonal sum of blocks of two
kinds:
\begin{align*}
 A(a,d)&:\quad q(e)=a,\quad b(e,e)=d\in R^\times,\\
 H(a_1,a_2;c)&:\quad
 b=\begin{pmatrix}0&c\\c&0\end{pmatrix},\quad c\in R^\times.
\end{align*}
A regular pair is basic metabolic exactly when it has no loop blocks and each
edge block has \(a_1=0\) or \(a_2=0\).  It is split exactly when every block is
\(H(0,0;c)\).
\end{theorem}

\begin{proof}
If \(b^\flat\) is invertible, its matrix sends a basis of \(V\) to a basis of
\(V^*\).  Unique bases make it a generalized permutation matrix.  Symmetry
makes the underlying permutation an involution.  A fixed point gives a loop
block; a transposition gives an edge block.

Let \(L\) be basic with \(L=L^\perp\).  A loop basis vector can neither lie in
\(L\), because it is not self-orthogonal, nor outside \(L\), because it would
then lie in \(L^\perp\setminus L\).  In an edge block, equality with the
orthogonal complement forces \(L\) to contain exactly one endpoint.  The
condition \(q|_L=0\) says that the selected endpoint has zero quadratic label.
Complementary Lagrangians exist exactly when both endpoint labels vanish.
\end{proof}

Consequently, whenever admissibility is decidable, the stable invariant is
decidable: decompose the graph, canonically compare the finitely many gauge
orbits, delete the metabolic connected edge types, and sort the remaining
multiset.

\section{The Hessenberg arm}

Let \(\#\) and \(\otimes\) denote Hessenberg natural sum and product.  Their
Cantor-normal-form definitions add coefficients and formally multiply
monomials, using natural sum on the exponents \cite{carruth1942}.

\begin{definition}[Closed finite-CNF fragment]
For \(d\in\Nat\), put
\[
 \Hess_d=\{\alpha:\alpha<\omega^{\omega^d}\}.
\]
It is equipped with \(\#\) and \(\otimes\), not ordinary ordinal addition or
multiplication and not Conway nim arithmetic.
\end{definition}

\begin{proposition}[Polynomial model]\label{prop:polynomial}
There is a semiring isomorphism
\[
 (\Hess_d,\#,\otimes)\cong\Nat[x_0,\ldots,x_{d-1}].
\]
Hence addition is cancellative, nonzero elements are closed under both
operations, \(\Hess_d^\times=\{1\}\), \(a\#a=1\) has no solution, and
\(a\#a=0\) implies \(a=0\).
\end{proposition}

\begin{proof}
Every exponent \(\beta<\omega^d\) has the unique expansion
\(\beta=\sum_{k<d}\omega^kn_k\).  Send
\(\omega^\beta\) to \(\prod_{k<d}x_k^{n_k}\), and extend by the finite outer
Cantor normal form.  Natural sum adds like coefficients, while natural product
adds exponent vectors and multiplies coefficients.  This proves the
isomorphism.  All listed properties are immediate in a polynomial semiring
over \(\Nat\).
\end{proof}

Fix a quadratic pair over \(\Hess_d\).  We continue to write addition and
multiplication in semiring notation.

\begin{lemma}[Unique balanced companion]\label{lem:balanced}
The companion is uniquely determined by \(q\), and
\[
 b(x,x)=2q(x)
\]
for every \(x\).
\end{lemma}

\begin{proof}
If \(b\) and \(c\) are companions, cancel \(q(x)+q(y)\) from their two
instances of the quadratic-pair identity.  For balance, compute \(q(x+x)\) in
two ways:
\[
 4q(x)=2q(x)+b(x,x).
\]
Additive cancellation gives \(b(x,x)=2q(x)\).
\end{proof}

Thus a rank-two pair has the unique form
\[
 Q(a,c,d)(x,y)=ax^2+cxy+dy^2,\qquad
 b=\begin{pmatrix}2a&c\\c&2d\end{pmatrix}.
\]

\begin{proposition}[Rank two]\label{prop:hessranktwo}
The isometry class of \(Q(a,c,d)\) is the orbit of \((a,c,d)\) under
\((a,c,d)\leftrightarrow(d,c,a)\).  It is decomposable exactly when \(c=0\),
and regular exactly when
\[
 (a,c,d)=(0,1,0).
\]
\end{proposition}

\begin{proof}
The only scalar unit is one, so every rank-two automorphism is either the
identity or the coordinate swap.  The graph is disconnected exactly when its
one possible edge is absent.  If the Gram matrix is invertible, unique bases
make it a permutation matrix.  Its diagonal entries \(2a,2d\) cannot equal
one.  Thus it must be the transposition matrix, which forces the displayed
triple.  The converse is direct.
\end{proof}

Write \(H=Q(0,1,0)\).

\begin{theorem}[Hessenberg regular collapse]\label{thm:hesscollapse}
Every regular finite free quadratic pair over \(\Hess_d\) is isometric to
\(H^{\perp m}\) for a unique \(m\).  Regular, split, and basic metabolic pairs
coincide.  Consequently
\[
 \operatorname{GW}_{\mathrm{reg}}(\Hess_d)\cong\Integers[H],
 \qquad W_{\mathrm{reg}}(\Hess_d)=0.
\]
\end{theorem}

\begin{proof}
Theorem~\ref{thm:blocks} reduces the issue to its permutation orbits.  A loop
would have \(1=b(e,e)=2q(e)\), impossible by
Proposition~\ref{prop:polynomial}.  Hence the permutation is fixed-point-free.
On a transposition the only nonzero Gram entries are the two off-diagonal
ones, both equal to the only unit, one.  Its diagonal equality
\(0=2q(e)\) forces both quadratic labels to vanish.  Every orbit is therefore
\(H\).  The remaining assertions follow from the block theorem.
\end{proof}

The singular envelope does not collapse.  Let \(\Connected_d\) be the set of
connected \(\Hess_d\)-weighted graph types from
Definition~\ref{def:graph}.

\begin{corollary}[Complete stable invariant]\label{cor:hessstable}
\[
 \Pair(\Hess_d)/{\sim_{\mathrm{st}}}
 \cong\Nat^{(\Connected_d\setminus\{H\})},
 \qquad
 W_{\mathrm{all}}(\Hess_d)
 \cong\Integers^{(\Connected_d\setminus\{H\})}.
\]
There are no further relations.
\end{corollary}

\begin{proof}
Theorem~\ref{thm:graph} gives the free monoid on all connected types, and
Theorem~\ref{thm:hesscollapse} says that the metabolic submonoid is freely
generated by the single type \(H\).  Quotienting and group-completing deletes
exactly that generator.
\end{proof}

For example, \(Q(0,2,0)\) and \(H\) have identical diagonal data, but the
former is a connected generator different from \(H\) and survives stably.
Thus the singular invariant contains genuine off-diagonal information beyond
the forced Cantor-normal-form labels.  An exact decision procedure normalizes
all CNFs, finds connected components, compares their finitely many vertex
permutations, discards \(H\), and compares the sorted multisets.

There is a useful warning about definitions.  If regularity is omitted and
``metabolic'' means only that some basic \(L\) satisfies \(q|_L=0\) and
\(L=L^\perp\), then every \(Q(0,c,0)\) with \(c\ne0\) becomes metabolic: take
\(L=\Hess_de_1\).  When \(c\) is not a unit the polar map is not surjective.
Such a definition would erase singular edge weights for an accidental
semiring reason, so regularity or perfect cross-pairing is indispensable.

\section{The rational supertropical arm}

Let \(U_{\Rationals}\) be the standard tangible max-plus supertropical cover
\[
\{0\}\sqcup\{t^\gamma:\gamma\in\Rationals\}
 \sqcup\{(t^\gamma)^\nu:\gamma\in\Rationals\}.
\]
Unequal values add by dominance; a tie becomes the corresponding ghost.  The
tangible elements \(t^\gamma\) are exactly the scalar units.  Write
\(e=1+1\) for the ghost of one.  The
value group is dense and two-divisible.  The min-plus statement is obtained by
\(\gamma\mapsto-\gamma\).  Free modules have unique bases
\cite[Theorem~0.9]{ikr2016qf1}.

For a pair, label vertex \(i\) by \((a_i,d_i)\) and edge \(ij\) by
\(c_{ij}\), as in Definition~\ref{def:graph}.  Theorem~\ref{thm:graph} gives
the following complete presentation.

\begin{theorem}[Supertropical stable blocks]\label{thm:superstable}
The pair monoid over \(U_{\Rationals}\) is free commutative on connected
admissible coefficient graphs modulo
\[
 (a_i,d_i,c_{ij})\mapsto
 (u_i^2a_{\pi(i)},u_i^2d_{\pi(i)},u_iu_jc_{\pi(i)\pi(j)}).
\]
Every regular pair is an orthogonal sum of loop blocks \(A(a,d)\) and edge
blocks \(H(a_1,a_2;c)\) from Theorem~\ref{thm:blocks}.  Its stable group is
free abelian on all connected types except the edge types
\[
 H(0,a;c),\qquad H(a,0;c),\qquad c\in U_{\Rationals}^\times,
\]
which are exactly the connected basic metabolic types.
\end{theorem}

This is a generators-and-relations theorem with no hidden relations.  It also
gives a decision procedure.  For each of finitely many permutations, write
\(u_i=t^{s_i}\).  Equality of the value labels becomes a finite rational
linear system in the \(s_i\), together with equality of zero patterns and of
tangible/ghost layers.  Canonical connected-orbit representatives can
therefore be compared in every finite rank.

Companion ambiguity is already visible in rank two.  Present a functional
quadratic form as
\[
 q=\begin{bmatrix}a_1&\alpha\\&a_2\end{bmatrix}.
\]

\begin{proposition}[Complete rank-two companions]\label{prop:supercompanion}
All companions have matrices
\[
 b=\begin{pmatrix}d_1&c\\c&d_2\end{pmatrix},
 \qquad d_i\in[0,ea_i],
\]
where
\[
 c\in
 \begin{cases}
 \{\alpha\},&a_1a_2=0,\\
 \{\alpha\},&\alpha^2>_\nu a_1a_2,\\
 \{z:z^2\le_\nu a_1a_2\},&a_1a_2\ne0\text{ and }
     \alpha^2\le_\nu a_1a_2.
 \end{cases}
\]
The choices of \(d_1,d_2,c\) are independent.
\end{proposition}

\begin{proof}
The diagonal intervals are
\cite[Proposition~6.12]{ikr2016qf1}.  If one diagonal coefficient vanishes,
the off-diagonal companion is rigid by Proposition~7.9 there.  Otherwise the
dense, square-divisible specialization of Theorems~7.11 and~7.12 gives the
remaining two cases.  Independence of the matrix entries is
Theorem~6.3 of the same source.
\end{proof}

In particular, one functional \(q\) can admit companions with different
regular block structures.  Stable theory over a supertropical semiring must
therefore specify pairs, not only functional quadratic forms.

\subsection{The scalar-extension obstruction}

Let \(K\) be a field of characteristic different from two and let
\(\varphi:K\to U_{\Rationals}\) be a tangible supervaluation into the standard
rational cover.  Thus \(\varphi(a)=t^{v(a)}\) for \(a\ne0\), where
\(v:K^\times\to\Rationals\), and the tangible unit group of
\(U_{\Rationals}\) is square-divisible.  Coefficientwise
supertropicalization is functorial for an ordered based form and monoidal under
orthogonal sum \cite[Definition~9.3]{ikr2016qf1}.  It is not an invariant of
the unbased form.

Consider the classical hyperbolic plane
\[
 q(xe+yf)=xy.
\]
In the basis \((e,f)\), its supertropicalization with the canonical balanced
companion is the split edge \(H(0,0;1)\), hence has stable class zero.  In the
equally valid basis
\[
 v_1=e+f,\qquad v_2=e-f,
\]
one has \(q(rv_1+sv_2)=r^2-s^2\).  Equations (9.11)--(9.13) of
\cite{ikr2016qf1} give
\[
 D(1,e)\perp D(1,e),
 \qquad D(a,ea): q(\epsilon)=a,\quad b(\epsilon,\epsilon)=ea.
\]
The ghost \(e\) is not a scalar unit, so both loop blocks are nonregular.
They are nevertheless connected generators in the all-pairs group and are
not metabolic.  Therefore
\[
 [H(0,0;1)]=0,
 \qquad 2[D(1,e)]\ne0.
\]

\begin{corollary}[No unbased canonical extension]\label{cor:nofunctor}
Canonical coefficientwise supertropicalization does not descend to unbased
isometry classes, even after regular basic-metabolic stabilization.  Moreover,
the regular category is not closed under changing the classical input basis.
\end{corollary}

Remark~9.5 of \cite{ikr2016qf1} supplies a distinct choice: tropicalize the
classical polar companion itself.  Its diagonal is \(\varphi(2a_i)\), rather
than \(e\varphi(a_i)\), and it is still a companion.  In the changed basis
above this alternative produces two regular nonmetabolic loops; it repairs
regularity in this example but not basis independence.  This distinction is
also the basis-dependence phenomenon studied explicitly in
\cite[Section~8]{ikr2018qf2}.

We now impose exactly the missing relations.  For
\(\epsilon\in\{\mathrm{bal},\mathrm{pol}\}\), let
\(S_\epsilon(K,\varphi)\) be the subgroup of
\(W_{\mathrm{all}}(U_{\Rationals})\) generated by the based lifts
\(T_\epsilon(q,E)\) of nonsingular classical forms.  Let
\(\overline S_\epsilon(K,\varphi)\) be its quotient by all relations
\[
 [T_\epsilon(q,E)]=[T_\epsilon(q,E')]
\]
for two bases of the same form, together with classical hyperbolic
stabilization.  This is the universal target for an additive
basis-independent invariant extracted from the indicated coefficientwise
lift.

\begin{theorem}[Universal ring-image collapse]\label{thm:ringcollapse}
For every \(K,\varphi\) as above,
\[
 \overline S_{\mathrm{bal}}(K,\varphi)\cong\Integers/2,
 \qquad
 \overline S_{\mathrm{pol}}(K,\varphi)\cong\Integers/2.
\]
Both isomorphisms are rank modulo two.  Consequently no square class,
discriminant, Hasse invariant, signature, Witt index, coefficient valuation,
or residue datum survives the universal unbased ring-image quotient.
\end{theorem}

\begin{proof}
Diagonalize a nonsingular form:
\[
 q\cong\langle a_1,\ldots,a_n\rangle,
 \qquad a_i\ne0.
\]
For each \(i\), choose a tangible unit \(u_i\) with
\(u_i^2=\varphi(a_i)^{-1}\).  In the canonical balanced lift, scaling the
\(i\)-th basis vector by \(u_i\) normalizes its rank-one pair to
\[
 D:\qquad q(\epsilon)=1,\quad \beta(\epsilon,\epsilon)=e.
\]
Hence every \(n\)-dimensional canonical image is basis-related to
\(D^{\perp n}\).

For the polar-companion lift, multiplicativity gives
\[
 \varphi(2a_i)=\varphi(2)\varphi(a_i).
\]
The same scaling normalizes every line to the single pair
\[
 A_{\varphi(2)}:\qquad
 q(\epsilon)=1,\quad b(\epsilon,\epsilon)=\varphi(2).
\]
Thus the polar image of every \(n\)-dimensional form is basis-related to
\(A_{\varphi(2)}^{\perp n}\).  In particular, valuations and square classes
change neither generator.

For the hyperbolic plane \(q(xe+yf)=xy\), the basis \((e,f)\) gives the split
edge \(H(0,0;1)\), which is zero in the stable group.  The basis
\((e+f,e-f)\) gives \(r^2-s^2\), hence two copies of \(D\), respectively two
copies of \(A_{\varphi(2)}\).  Therefore the relevant generator \(L\) obeys
\(2L=0\), proving that each quotient is at most \(\Integers/2\).

Rank modulo two is additive, unchanged by basis change, and zero on every
basic metabolic pair.  It sends \(L\) to one, so \(L\ne0\) and no stronger
relation is possible.  This also proves the stated universal property.
\end{proof}

The theorem concerns the subgroup generated by classical scalar-extension
images.  It does not collapse the intrinsic group
\(W_{\mathrm{all}}(U_{\Rationals})\): connected noncanonical weight or layer
types that no nonsingular classical lift realizes remain independent free
generators.  In particular, the isolated zero line survives in the ambient
group even though it is absent from the nonsingular ring-image subgroup.

There is also a sharp regularity distinction.  The polar lift has a regular
diagonal presentation, and its regular ring-image quotient is again
\(\Integers/2\).  A generic polar presentation is nevertheless nonregular.
For the canonical balanced lift, a nonzero diagonal companion coefficient is
ghost and not a unit.  A regular canonical ring image must therefore have
zero diagonal and perfect-matching off-diagonal support, which makes the
classical form hyperbolic and the image split.  Thus its regular ring-image
subgroup is zero; the all-pairs category in
Theorem~\ref{thm:ringcollapse} is essential.

If degenerate classical forms are included, the standard cover still leaves
only total-rank parity.  Indeed, the polar matrices
\[
 \begin{pmatrix}1&1\\1&1\end{pmatrix},
 \qquad
 \begin{pmatrix}1&1\\1&-1\end{pmatrix}
\]
have identical coefficientwise images, while their ranks are one and two.
After classical diagonalization this identifies the zero line \(Z\) with
\(L\); hence every finite free form has class \(\rank(V)L\).

\begin{proposition}[Trivially valued characteristic two]\label{prop:char2collapse}
Let \(K\) have characteristic two and let \(\varphi(0)=0\),
\(\varphi(K^\times)=1\).  For the canonical balanced and polar-companion lifts
considered separately, the nonsingular universal ring-image quotient is
\[
 \begin{cases}
  \Integers/2,&K=\mathbb F_2,\\
  0,&|K|>2.
 \end{cases}
\]
The surviving \(\mathbb F_2\)-class is the Arf invariant.
\end{proposition}

\begin{proof}
The polar form is nondegenerate alternating, so a symplectic basis decomposes
\(q\) into planes
\[
 a_ix^2+xy+b_iy^2.
\]
Only the zero patterns of \((a_i,b_i)\) remain after \(\varphi\).  The
\((0,0)\)-plane is split.  A plane with exactly one nonzero endpoint is also
killed: for the balanced lift this follows by comparing the bases
\((e,f)\) and \((e,e+f)\) of \(q=xy\), while for the polar lift it is already
basic metabolic.

It remains only the plane with both endpoints nonzero.  If \(|K|>2\), choose
\(\lambda\notin\{0,1\}\).  The hyperbolic basis
\[
 e+f,\qquad e+\lambda f
\]
has two nonzero quadratic labels and nonzero mutual polar label, so the
remaining generator is itself zero.  Over \(\mathbb F_2\), that generator is
the anisotropic plane \(x^2+xy+y^2\).  Its coefficient support determines the
Arf bit \cite{arf1941}.  For two copies with bases
\((e_1,f_1,e_2,f_2)\), the vectors
\[
 f_1+f_2,\quad f_1+e_2,\quad
 e_1+e_2+f_2,\quad e_1+f_1+e_2+f_2
\]
are isotropic and form two mutually orthogonal symplectic pairs.  Thus two
anisotropic planes are hyperbolic, giving order two, while Arf proves that one
does not vanish.  Symplectic decomposition proves the assertion in every
rank.
\end{proof}

\section{Thermograph walls}

Let \(\Gamma\) be the additive group of continuous functions
\(f:[0,\infty)\to\mathbb R\) that are piecewise affine with finitely many
rational breakpoints, rational slopes and intercepts, and are eventually
constant.  Pointwise maximum, minimum, negation, and division by two preserve
\(\Gamma\).  Define
\[
 \mathbb P_{\max}=\Gamma\sqcup\{-\infty\},\qquad
 f\oplus g=\max(f,g),\quad f\odot g=f+g.
\]
This is an entire zerosumfree idempotent semifield.  Its finite elements are
all units and its square map is surjective.  Negation gives the dual
min-plus semifield \(\mathbb P_{\min}=\Gamma\sqcup\{+\infty\}\).

Let \(L,R\in\Gamma\).  On the free rank-one
\(\mathbb P_{\max}\)-module define
\[
 q_L(x)=x^{\odot2}\odot L,
 \qquad b_R(x,y)=x\odot y\odot R.
\]

\begin{theorem}[Wall pair and gap classification]\label{thm:wallpair}
The pair \((q_L,b_R)\) is quadratic exactly when \(R\le L\) pointwise.  It is
regular, and
\[
 Q(L,R)\cong Q(L',R')
 \quad\Longleftrightarrow\quad
 L-R=L'-R'.
\]
\end{theorem}

\begin{proof}
Homogeneity and bilinearity are immediate.  Pointwise,
\begin{align*}
 q_L(x\oplus y)&=\max(2x+L,2y+L),\\
 q_L(x)\oplus q_L(y)\oplus b_R(x,y)
 &=\max(2x+L,2y+L,x+y+R).
\end{align*}
If \(R\le L\), the third term is at most the maximum of the first two.  Taking
\(x=y\) equal to the tropical multiplicative unit proves the converse.

The Gram coefficient \(R\) is finite and hence a unit, proving regularity.
Every rank-one automorphism is multiplication by a finite unit \(u\), which is
an isometry precisely when
\[
 L=L'+2u,\qquad R=R'+2u.
\]
This preserves the gap.  Conversely, equal gaps give such a unit with
\(u=(L-L')/2\).
\end{proof}

The left and right thermograph walls of a short game satisfy \(R_G\le L_G\)
and belong to \(\Gamma\) \cite{siegel2013}.  Thus they define \(Q(G)=
Q(L_G,R_G)\), whose isometry class is exactly the gap \(D_G=L_G-R_G\).
The dual min-plus construction records the same gap after negation.

In the stable group of Definition~\ref{def:stable}, distinct rank-one loop
types survive freely: a basic metabolic regular object has only edge blocks by
Theorem~\ref{thm:blocks}.  The necessary basepoint reduction is
\[
 \Omega(G)=[Q(G)]-[Q(0)].
\]

\begin{theorem}[Hot-switch obstruction]\label{thm:wallnogo}
The map \(\Omega\) is not additive under disjunctive sum.  Consequently the
paired max/min wall class is not an additive stable game invariant.
\end{theorem}

\begin{proof}
Take the switch \(X=\{1\mid-1\}\).  Its walls and gap are
\[
 L_X(t)=\max(1-t,0),\qquad
 R_X(t)=-\max(1-t,0),\qquad
 D_X(t)=2\max(1-t,0).
\]
The zero game has gap zero.  Theorem~\ref{thm:wallpair} and freeness of the
rank-one generators imply that \(\Omega(X)\ne0\) and has infinite order.

But \(-X=X\), so \(X+X=0\) in the Conway game group
\cite{conway1976}.  Additivity would give
\(2\Omega(X)=\Omega(0)=0\), a contradiction.  Projection of a paired max/min
invariant to its max component would be additive, so the same contradiction
rules out the pair.  A synchronized quotient can evade it only by killing the
nonzero gap atom exhibited above.
\end{proof}

There is a complementary freezing obstruction at temperature zero.  Any
additive invariant into a cancellative target that depends only on the frozen
thermograph assigns the same value to \(\uparrow\) and \(2\uparrow\).  The
identity \(x=2x\) then cancels to \(x=0\).  Thus wall-only additivity kills the
entire common infinitesimal wall class even before the hot-switch argument.

\section{Formal boundary and implementation boundary}

The Lean module \texttt{Ogdoad.SemiringQuadratic} \cite{ogdoad2026} checks:
\begin{enumerate}
\item uniqueness and diagonal balance of companions over additively
cancellative semirings;
\item that finite basis labels and Gram coefficients determine the entire
quadratic pair;
\item the no-half-of-one, no-nonzero-self-double, and trivial-unit properties
of the polynomial model
\[
 \operatorname{MvPolynomial}(\operatorname{Fin}d,\Nat)
\]
for \(\Hess_d\);
\item that every two-sided invertible square matrix over that polynomial
semiring is a permutation matrix;
\item the end-to-end polynomial proxy theorem
\texttt{hessenbergPolynomial\_}\linebreak
\texttt{regular\_isometric\_}\linebreak
\texttt{hyperbolic}: bijectivity of the
companion adjoint implies even rank and produces an explicit coefficient
calculation together with an actual linear isometry of quadratic pairs to a
canonical orthogonal sum of hyperbolic planes; and
\item that an additive invariant into a cancellative monoid vanishes on any
observation identified with its double; and
\item Mathlib's classical diagonalization under \(2\ne0\), together with an
explicit additive equivalence between the exact cyclic/hyperbolic/parity
presentation of the ring-image group and \(\Integers/2\).
\end{enumerate}
Thus the kernel now checks the regular hyperbolic classification over the
concrete polynomial model from the same input notion of regularity used in the
paper; it no longer assumes the permutation-Gram, matching, or
coefficient-to-isometry joins.  The
Cantor-normal-form identification with the actual ordinal fragment, the
connected-graph stable quotient and Witt-group packaging, the construction of
the concrete supertropical semiring and its two scalar-extension presentations,
and standard thermograph facts remain proved here or cited external
mathematics, not Lean theorems.  For Theorem~\ref{thm:ringcollapse}, Lean
checks classical diagonalization and the exact final group equivalence, while
the supertropical normalization and quotient-presentation fields are the
explicit remaining bridge.

Ogdoad's current \texttt{Ordinal} backend implements Conway nim arithmetic
and ordinary noncommutative Cantor arithmetic.  It is not a Hessenberg
semiring.  Likewise, the current finite-valued piecewise-linear wall type has
no global tropical zero.  Reusing either type silently would violate the
coefficient hypotheses above.  A future implementation should add a distinct
finite-CNF natural-arithmetic type and the explicit one-point completions
\(\mathbb P_{\max},\mathbb P_{\min}\).

\section{Conclusion}

Unique bases turn stable quadratic-pair theory into connected weighted graph
theory.  In the Hessenberg fragments this gives both extremes at once: the
regular Witt theory collapses completely, while the singular stable envelope
is a large free invariant with a finite decision procedure.  Over the rational
supertropical cover the same graph theorem gives exact blocks, companions,
and metabolic relations.  Coefficientwise supervaluation freezes the chosen
classical basis, but the universal quotient imposing all source basis changes
is now explicit: for both standard companion lifts it is exactly rank parity
in characteristic different from two.

Thermograph walls do form honest rank-one quadratic pairs after the minimal
tropical completion, and their stable class is exactly the wall gap.  The
self-negative hot switch proves that this geometrically natural class is
incompatible with additivity.  Thus the three coefficient-world questions
posed here are closed: Hessenberg regular collapse with a nontrivial singular
envelope, exact supertropical stable blocks with the universal ring-image
quotient, and the thermograph wall obstruction.

\bibliographystyle{plain}
\bibliography{semiring_stability}

\end{document}
